\section{Results --- the atlas}
\label{sec:results}

The tables and figures in this section are generated from the campaign by
\texttt{benchmarks/analyze.py} (Sec.~\ref{sec:repro}); every number and every reconstruction is
reproducible from the run it names. \texttt{\textbackslash IfFileExists} guards let the document
build before the run has finished.

\subsection{Which solver wins, and where}
For each (structure, noise) the solver with the best NMSE, and its curated metrics.
\IfFileExists{tables/cross_comparison_matirf.tex}{\begin{center}\begin{minipage}{\linewidth}\centering
\captionof{table}{Best algorithm by NMSE on MA-TIRF --- winner per structure and noise level.}
\begin{adjustbox}{max width=\linewidth}
\begin{tabular}{lllllllll}
\toprule
dataset & noise & best (by NMSE) & params & NMSE & PSNR & Depth Error (nm) & Stack Recovery & chi2\_ratio \\
\midrule
cell & 0 & \textbf{ADMM} & kappa=0.1, mu=0.5 & 0.0318 & 28.8 & 4 & 0.636 & 5.2e-08 \\
cell & 0.02 & \textbf{ADMM-PnP} & denoiser=Gaussian, sigma=50, rho=0.03 & 0.177 & 21.3 & 9 & 0 & 0.96 \\
cell & 0.05 & \textbf{ADMM-PnP} & denoiser=Gaussian, sigma=50, rho=0.03 & 0.307 & 18.9 & 21 & 0 & 0.903 \\
fibres & 0 & \textbf{ADMM} & kappa=0.01, mu=0.5 & 0.0406 & 38.0 & 10 & 0.966 & 4.86e-06 \\
fibres & 0.02 & \textbf{PNP} & denoiser=Gaussian, sigma=25, lambda\_kz=12, iter=40 & 0.385 & 28.2 & 59 & 0.0172 & 0.527 \\
fibres & 0.05 & \textbf{PNP} & denoiser=Gaussian, sigma=25, lambda\_kz=12, iter=40 & 0.653 & 25.9 & 85 & 0.069 & 0.459 \\
vesicles & 0 & \textbf{ADMM} & kappa=0.01, mu=0.5 & 0.021 & 42.9 & 8 & 1 & 5.3e-06 \\
vesicles & 0.02 & \textbf{MCMC} & denoiser=TV Bregman, sigma=0.02, lambda\_rr=38.863, iter=300 & 0.512 & 29.0 & 61 & 1 & 0.609 \\
vesicles & 0.05 & \textbf{PNP} & denoiser=Gaussian, sigma=25, lambda\_kz=12, iter=40 & 0.732 & 27.5 & 98 & 0 & 0.435 \\
\bottomrule
\end{tabular}\end{adjustbox}\end{minipage}\end{center}
}%
  {\emph{[pending: run \texttt{python -m benchmarks analyze -{}-minimal}]}}
\IfFileExists{tables/cross_comparison_deconv.tex}{\begin{center}\begin{adjustbox}{max width=\linewidth}
\begin{tabular}{llllllll}
\toprule
dataset & noise & best (by NMSE) & params & NMSE & PSNR & SSIM & chi2\_ratio \\
\midrule
img\_001 & 0 & \textbf{ADAM} & reg=none, lambda=0 & 0.0191 & 24.4 & 0.694 & 1.96e-07 \\
img\_001 & 0.02 & \textbf{MCMC} & denoiser=TV Bregman, sigma=0.03 & 0.0245 & 23.4 & 0.596 & 0.98 \\
img\_001 & 0.05 & \textbf{ADAM} & reg=tikhonov, lambda=0.1 & 0.0371 & 21.5 & 0.511 & 1.09 \\
\bottomrule
\end{tabular}\end{adjustbox}\end{center}
}{}

\subsection{ADAM in depth: the prior sets the ceiling}
\label{sec:adam-depth}
The cross-comparison above runs each solver with a single default prior. ADAM, however, sits
well below its own ceiling until the prior is matched to the structure. To show this we swept the
regularisation weight $\lambda$ for ADAM on the reference truth \texttt{cell\_fibres\_vesicles}
at noise $0.02$ --- the same object and noise as the headline row above --- with the fully
converged settings used throughout the atlas (\texttt{lr}$=0.1$, ridge start, ridge weight
$s_2^2 = 36.24$, \texttt{EPS}$=10^{-8}$, $5000$ iterations). This refines the campaign's coarse
$\lambda$ grid, of which it is a superset.

\paragraph{Sparsity (L1) has a clear optimum; total variation does not.}
\begin{center}\begin{adjustbox}{max width=\linewidth}
\begin{tabular}{lcc}
\toprule
\multicolumn{3}{c}{\textbf{ADAM + L1}} \\
$\lambda$ & NMSE & $z$-err (nm) \\
\midrule
0.01 & 0.618 & 50 \\
0.02 & 0.590 & 45 \\
0.05 & 0.512 & 34 \\
0.10 & 0.401 & 25 \\
0.15 & 0.319 & 20 \\
0.20 & 0.270 & 16 \\
\textbf{0.30} & \textbf{0.268} & \textbf{14} \\
0.50 & 0.477 & 22 \\
\bottomrule
\end{tabular}
\hskip 3em
\begin{tabular}{lcc}
\toprule
\multicolumn{3}{c}{\textbf{ADAM + TV}} \\
$\lambda$ & NMSE & $z$-err (nm) \\
\midrule
\textbf{0.005} & \textbf{0.480} & 42 \\
0.010 & 0.487 & 35 \\
0.015 & 0.505 & 36 \\
0.020 & 0.519 & 38 \\
0.030 & 0.542 & 41 \\
0.050 & 0.571 & 40 \\
0.070 & 0.589 & 41 \\
0.100 & 0.614 & 43 \\
\bottomrule
\end{tabular}
\end{adjustbox}\end{center}

\noindent L1 has a well-defined minimum at $\lambda \approx 0.2$--$0.3$ (NMSE $0.27$, depth error
$14$~nm). TV, on the same object, is \emph{monotone increasing} in $\lambda$: its penalty only
erodes the depth structure it is meant to regularise, and its floor ($\approx 0.48$) stays far
above L1's. On a composite of filaments and vesicles the sparsity prior is the right one and
total-variation smoothing is not. Tuned this way ADAM closes most of the gap to the ADMM-PnP
winner (NMSE $0.17$) at a fraction of its runtime.

\paragraph{Across the atlas, the best ADAM prior tracks the structure.}
The setting that minimised NMSE for ADAM at each (structure, noise), over the full prior grid
(none / L1 / TV / SHV):
\begin{center}\begin{adjustbox}{max width=\linewidth}
\begin{tabular}{lllll}
\toprule
structure & noise & best prior & NMSE & $z$-err (nm) \\
\midrule
\texttt{cell\_fibres\_vesicles} & 0.02 & L1, $\lambda=0.2$ & 0.27 & 16 \\
\texttt{cell\_fibres\_vesicles} & 0.05 & L1, $\lambda=0.1$ & 0.62 & 29 \\
\texttt{vesicles} & 0 & L1, $\lambda=0.02$ & 0.068 & 6 \\
\texttt{vesicles} & 0.02 & TV, $\lambda=0.05$ & 0.30 & 46 \\
\texttt{vesicles} & 0.05 & TV, $\lambda=0.02$ & 0.69 & 102 \\
\texttt{deconv} (2D) & 0.02 & SHV, $\lambda=0.5$ & 0.026 & --- \\
\texttt{deconv} (2D) & 0.05 & SHV, $\lambda=0.5$ & 0.033 & --- \\
\bottomrule
\end{tabular}\end{adjustbox}\end{center}

\noindent Sparsity for point-like and composite structures, curvature (SHV) for the smooth 2D
deconvolution target: the winning prior follows the object. This is the report's thesis in one
table --- the prior, not the optimiser, sets the ceiling.

\paragraph{PPXA is the costly, fragile member here.} On the same problem PPXA reaches at best
NMSE $0.46$ (vesicles, noise $0.02$) but at a median $610$~s per run against ADAM's $120$~s, with
most of its configurations rejected by the realism check or timed out. Where ADAM's prox-free
gradient scheme converges cheaply, the proximal splitting buys no accuracy at several times the
cost.

\subsection{Reconstructions by structure and noise}
The visual heart of the report: each solver's best reconstruction (depth map + $yz$/$zx$
profiles, exported from the GUI viewers). Figures are generated into \texttt{figures/} by the
analysis; a fuller grid can be assembled from the files there.

\begin{figure}[H]\centering
  \includegraphics[width=.9\linewidth]{matirf_cell_fibres_vesicles_ADMM_noise00}
  \caption{\texttt{cell\_fibres\_vesicles}, no noise --- \textbf{ADMM} (the winner: NMSE 0.04,
  depth error 4~nm). At zero noise the sparsity prior localises the structure almost exactly.}
\end{figure}

\begin{figure}[H]\centering
  \includegraphics[width=.49\linewidth]{matirf_cell_fibres_vesicles_ADMM-PnP_noise005}\hfill
  \includegraphics[width=.49\linewidth]{matirf_cell_fibres_vesicles_PNP_noise005}
  \caption{\texttt{cell\_fibres\_vesicles}, noise 0.05. \emph{Left:} \textbf{ADMM-PnP}, the
  best under noise (NMSE 0.37) --- the denoiser prior holds the structure. \emph{Right:}
  \textbf{PnP-HQS}, which collapses at the same noise. Same object, same noise: the difference
  is the method's robustness.}
\end{figure}

\subsection{Effective parameter ranges}
The setting that minimised NMSE at each noise level --- the atlas proper. Compare to the
predicted ranges in \texttt{docs/algorithms/}.
\IfFileExists{tables/effective_ranges.tex}{\paragraph{ADAM}
\begin{center}\begin{adjustbox}{max width=\linewidth}
\begin{tabular}{llll}
\toprule
noise & best params & NMSE & depth err (nm) \\
\midrule
0 & reg=none, lambda=0 & 0.635 & 68 \\
0.02 & reg=tv, lambda=0.05 & 0.561 & 40 \\
0.05 & reg=tv, lambda=0.05 & 0.668 & 58 \\
\bottomrule
\end{tabular}\end{adjustbox}\end{center}
\paragraph{ADMM}
\begin{center}\begin{adjustbox}{max width=\linewidth}
\begin{tabular}{llll}
\toprule
noise & best params & NMSE & depth err (nm) \\
\midrule
0 & kappa=0.1, mu=0.5 & 0.0409 & 4 \\
0.02 & kappa=0.3, mu=1 & 0.608 & 28 \\
0.05 & kappa=0.3, mu=1 & 0.911 & 60 \\
\bottomrule
\end{tabular}\end{adjustbox}\end{center}
\paragraph{ADMM-PnP}
\begin{center}\begin{adjustbox}{max width=\linewidth}
\begin{tabular}{llll}
\toprule
noise & best params & NMSE & depth err (nm) \\
\midrule
0 & denoiser=Bilateral, sigma=25, rho=0.1 & 0.0524 & 4 \\
0.02 & denoiser=Gaussian, sigma=25, rho=0.1 & 0.174 & 13 \\
0.05 & denoiser=Gaussian, sigma=50, rho=0.1 & 0.367 & 26 \\
\bottomrule
\end{tabular}\end{adjustbox}\end{center}
\paragraph{MCMC}
\begin{center}\begin{adjustbox}{max width=\linewidth}
\begin{tabular}{llll}
\toprule
noise & best params & NMSE & depth err (nm) \\
\midrule
0 & denoiser=TV Bregman, sigma=0.01 & 0.406 & 19 \\
0.02 & denoiser=TV Bregman, sigma=0.01 & 0.432 & 21 \\
0.05 & denoiser=TV Bregman, sigma=0.01 & 0.588 & 38 \\
\bottomrule
\end{tabular}\end{adjustbox}\end{center}
\paragraph{PNP}
\begin{center}\begin{adjustbox}{max width=\linewidth}
\begin{tabular}{llll}
\toprule
noise & best params & NMSE & depth err (nm) \\
\midrule
0 & denoiser=Gaussian, sigma=50 & 0.462 & 27 \\
0.02 & denoiser=None & 0.887 & 81 \\
0.05 & denoiser=None & 0.936 & 83 \\
\bottomrule
\end{tabular}\end{adjustbox}\end{center}
\paragraph{PPXA}
\begin{center}\begin{adjustbox}{max width=\linewidth}
\begin{tabular}{llll}
\toprule
noise & best params & NMSE & depth err (nm) \\
\midrule
0.02 & reg=tv, lambda=0.1 & 0.619 & 40 \\
\bottomrule
\end{tabular}\end{adjustbox}\end{center}
}%
  {\emph{[pending the run]}}

\subsection{Rejections, failures and cost}
Rejection rate, dominant realism-check findings, runtime and peak memory per solver --- including
the PnP-family out-of-memory behaviour on the real \texttt{esoubies} stack.
\IfFileExists{tables/rejects_and_cost.tex}{\begin{center}\begin{minipage}{\linewidth}\centering
\captionof{table}{Cost and robustness per solver --- runs, rejection rate, failures, median runtime, peak memory and top findings.}
\begin{adjustbox}{max width=\linewidth}
\begin{tabular}{lllllll}
\toprule
solver & runs & rejected & failed/timeout & median s & max mem (MB) & top findings \\
\midrule
ADAM & 45 & 4 (8\%) & 0 & 17.2 & 2176 & no-better×72, depth-shift×18, collapsed×4 \\
ADMM & 19 & 1 (5\%) & 0 & 1.7 & 2204 & no-better×20, depth-shift×10, collapsed×1 \\
ADMM-PnP & 37 & 3 (8\%) & 1 & 14.7 & 10210 & no-better×61, depth-shift×18, collapsed×2 \\
MCMC & 39 & 0 (0\%) & 0 & 9.7 & 2236 & no-better×47, depth-shift×17 \\
PNP & 37 & 0 (0\%) & 0 & 3.7 & 3591 & no-better×40, depth-shift×18 \\
PPXA & 2 & 0 (0\%) & 0 & 291.7 & 1656 & no-better×3, depth-shift×1 \\
\bottomrule
\end{tabular}\end{adjustbox}\end{minipage}\end{center}
}%
  {\emph{[pending the run]}}

\subsection{Real data (\texttt{esoubies})}
No truth, so the reconstruction is judged by \texttt{chi2\_ratio} and by eye.

\begin{figure}[H]\centering
  \includegraphics[width=.9\linewidth]{matirf_esoubies_ADMM-PnP_native}
  \caption{\texttt{esoubies} (real measurement, no ground truth) --- reconstructed with
  ADMM-PnP. The PnP family reconstructs it but is the memory-hungry one on this full-size
  volume (Sec.~\ref{sec:limits}).}
\end{figure}
